\documentclass[11pt]{article}
\usepackage[margin=1in]{geometry}
\usepackage[T1]{fontenc}
\usepackage{graphicx}
\usepackage{hyperref}
\usepackage{enumitem}
\setlength{\parskip}{0.5em}
\setlength{\parindent}{0pt}

\title{diffusion models for video generation}
\author{arXiv GraphRAG --- phase-6 deep research}
\date{\today}

\begin{document}
\maketitle

\begin{abstract}
An outstanding, frontier-level review that masterfully connects historical GAN/VAE methods to the current state-of-the-art in Diffusion Transformers and world models.
\end{abstract}

\section{From Pixels to Temporal Coherence: A Survey of Diffusion Models for Video Generation}

The ability to synthesize coherent, high-fidelity video from textual or visual prompts represents one of the most significant challenges in modern artificial intelligence. While generative modeling has achieved remarkable success in the static image domain, video generation introduces a complex temporal dimension that requires maintaining both spatial structure and temporal continuity across time. This task is critical for a wide array of applications, ranging from digital entertainment and cinematic production to autonomous vehicle simulation and medical imaging, where the ability to model dynamic, realistic environments is paramount.

This review traces the technological evolution of the field, beginning with the foundations of generative modeling and the early dominance of Generative Adversarial Networks (GANs) and autoregressive models. While these early approaches laid the groundwork for sequential data modeling, they often struggled with long-term stability and high-resolution synthesis. The field underwent a fundamental paradigm shift with the transition to latent video diffusion and the emergence of robust text-to-video frameworks, which leveraged the denoising capabilities of diffusion processes to achieve unprecedented levels of visual quality and semantic alignment.

As the field matured, research moved beyond simple generation toward fine-grained controllability, specifically regarding motion, camera trajectories, and compositional elements. This progression led to the integration of Transformers to address scaling and efficiency, enabling the generation of increasingly long-form video content. However, as the complexity of these models grows, the necessity for rigorous evaluation, benchmarking, and the enforcement of physical consistency becomes more acute. This review concludes by examining the current frontier, characterized by the development of advanced foundation models and the pursuit of real-time synthesis.

\subsection{Foundations of Generative Modeling and Diffusion}

The evolution of generative modeling has transitioned from variational inference and adversarial training toward the current dominance of diffusion-based frameworks. Early advancements focused on stochastic variational inference, where the reparameterization trick enabled the optimization of variational lower bounds via standard stochastic gradient methods (\texttt{arXiv:1312.6114}). To mitigate issues such as posterior collapse, the VQ-VAE introduced discrete latent representations (\texttt{arXiv:1711.00937}). Parallel to these approaches, Generative Adversarial Networks (GANs) established a minimax framework for the simultaneous training of a generator and discriminator (\texttt{arXiv:1406.2661}), bypassing the need for Markov chains or unrolled approximate inference. Later, style-based architectures enhanced GANs by providing scale-specific control and improved disentanglement of latent factors (\texttt{arXiv:1812.04948}).

A paradigm shift occurred with the introduction of models inspired by non-equilibrium thermodynamics, which utilize iterative forward diffusion to destroy structure and learned reverse processes to restore it (\texttt{arXiv:1503.03585}). This culminated in Denoising Diffusion Probabilistic Models (DDPM), which established the mathematical equivalence between diffusion, denoising score matching, and Langevin dynamics using an \$\textbackslash{}epsilon\$-prediction parameterization (\texttt{arXiv:2006.11239}). While DDPM achieved state-of-the-art image synthesis, it was fundamentally limited to static image generation, offering no methodology or framework for temporal modeling (\texttt{arXiv:2006.11239}). This framework was later unified through stochastic differential equations (SDEs), which connected score-based modeling and diffusion into a single framework (\texttt{arXiv:2011.13456}).

Efficiency remains a critical differentiator in these models. While DDPMs are computationally intensive, Denoising Diffusion Implicit Models (DDIM) introduced non-Markovian processes that accelerate sampling by \$10\textbackslash{}times\$ to \$50\textbackslash{}times\$ in wall-clock time without altering the training objective (\texttt{arXiv:2010.02502}). As diffusion models demonstrated the ability to outperform GANs through classifier guidance (\texttt{arXiv:2105.05233}), the field moved toward latent-space operations. Latent Diffusion Models (LDM) decoupled perceptual compression from semantic generation by training in a low-dimensional latent space of pretrained autoencoders (\texttt{arXiv:2112.10752}). This approach significantly reduces computational requirements and enables flexible conditioning via cross-attention; however, it faces a trade-off where excessive downsampling can cause information loss and stagnating fidelity (\texttt{arXiv:2112.10752}). These foundations—ranging from the convolutional context-capturing of U-Nets (\texttt{arXiv:1505.04597}) to the attention-based parallelization of Transformers (\texttt{arXiv:1706.03762})—provide the structural and mathematical basis for modern video generation.

\subsection{Early Video Generation: GANs and Autoregressive Models}

Early research in video generation focused heavily on Generative Adversarial Networks (GANs), primarily through the lens of decomposing video into constituent components or optimizing temporal stability. Initial approaches, such as TGAN, utilized a dual-generator architecture—comprising a temporal generator and an image generator—to produce sequences from a single latent variable, employing singular value clipping to ensure stable end-to-end training [\texttt{arXiv:1611.06624}]. This concept of decomposition was significantly advanced by MoCoGAN, which unsupervisedly separated motion and content by utilizing a sequence of random vectors [\texttt{arXiv:1707.04993}]. While MoCoGAN provided a foundation for disentanglement, subsequent work sought to address the inherent limitations of such models. For instance, researchers reformulated video generation as a Markov Decision Process (MDP) to mitigate temporal inconsistencies like video freezing or looping, offering a more memory-efficient method for modeling long-term motion [\texttt{arXiv:1909.12400}].

As the field progressed, the focus shifted toward scaling and efficiency. DVD-GAN introduced a dual video discriminator architecture to scale GANs to high-complexity datasets like Kinetics-600, achieving high fidelity through a computationally efficient discriminator decomposition [\texttt{arXiv:1907.06571}]. Similarly, Temporal Shift GAN sought to avoid the high computational costs of 3D architectures by facilitating information exchange between neighboring time points to maintain spatio-temporal consistency [\texttt{arXiv:2004.01823}]. However, this approach faced criticism regarding its effectiveness compared to other methods [\texttt{arXiv:2107.07224}].

A major trend emerged in leveraging pre-trained image models to enhance video synthesis. Rather than training from scratch, some frameworks utilized a fixed, pre-trained image generator and a motion generator to discover latent space trajectories, achieving an order of magnitude increase in computational efficiency [\texttt{arXiv:2104.15069}]. StyleVideoGAN further refined this by separating the spatial domain (via a pre-trained StyleGAN) from the temporal domain, significantly reducing the required training data [\texttt{arXiv:2107.07224}]. This trajectory culminated in StyleGAN-V, which treats video as a continuous signal rather than discrete frames, using positional embeddings to enable the generation of arbitrarily long videos at high frame rates [\texttt{arXiv:2112.14683}].

The evolution of these models also necessitated better evaluation. The lack of qualitative metrics for temporal coherence and diversity led to the proposal of the Fréchet Video Distance (FVD) metric and the StarCraft 2 Videos benchmark [\texttt{arXiv:1812.01717}]. Recent advancements have moved toward implicit representations, such as DIGAN, which uses Implicit Neural Representations (INRs) and a motion discriminator to improve motion dynamics without requiring the observation of entire long sequences [\texttt{arXiv:2202.10571}].

\subsection{Transition to Latent Video Diffusion and Text-to-Video}

The evolution of video generation has transitioned from direct pixel-space modeling to more efficient latent-space representations and the repurposing of pre-trained text-to-image (T2I) architectures. Early approaches, such as Video Diffusion Models (VDM), focused on extending standard image diffusion architectures to video by enabling joint training from image and video data to reduce minibatch gradient variance and speed up optimization [\texttt{arXiv:2204.03458}]. While VDM introduced conditional sampling for spatial and temporal extension [\texttt{arXiv:2204.03458}], it faced significant computational hurdles. In contrast, MagicVideo addressed these efficiency concerns by utilizing a Latent Diffusion Model (LDM) with a 3D U-Net and a VideoVAE, achieving approximately 64x fewer computations (FLOPs) than VDM [\texttt{arXiv:2211.11018}]. Similarly, research into high-fidelity long video generation has leveraged low-dimensional 3D latent spaces and hierarchical diffusion to outperform pixel-space models under limited computational budgets, enabling the generation of sequences exceeding 1,000 frames [\texttt{arXiv:2211.13221}].

A significant trend in the literature is the leveraging of large-scale T2I models to facilitate text-to-video (T2V) synthesis. It has been demonstrated that scaling large pre-trained text-only language models, such as T5, improves image-text alignment and sample fidelity more effectively than scaling the image diffusion model itself [\texttt{arXiv:2205.11487}]. This insight underpins several methodologies for video generation. Some approaches focus on high-definition output through cascaded architectures; for instance, Imagen Video employs a cascade of spatial and temporal super-resolution models [\texttt{arXiv:2210.02303}], while VideoGen utilizes a cascaded latent diffusion module followed by flow-based temporal upsampling and an enhanced video decoder [\texttt{arXiv:2309.00398}].

Other researchers have sought to adapt existing T2I models through tuning or zero-shot repurposing. Tune-A-Video introduced a one-shot tuning strategy using a tailored spatio-temporal attention mechanism to avoid large-scale video dataset training [\texttt{arXiv:2212.11565}]. This contrasts with the zero-shot approach of Text2Video-Zero, which requires no additional training or optimization by reprogramming frame-level self-attention and enriching latent codes with motion dynamics [\texttt{arXiv:2303.13439}]. Moving beyond tuning, AnimateDiff offers a plug-and-play motion module that can animate personalized T2I models without model-specific tuning [\texttt{arXiv:2307.04725}]. More recently, Stable Video Diffusion (SVD) has scaled latent video diffusion by inserting temporal layers into a pre-trained 2D architecture, utilizing a systematic three-stage training strategy—text-to-image pretraining, video pretraining, and high-quality video finetuning—to achieve high-fidelity results [\texttt{arXiv:2311.15127}]. Finally, ModelScopeT2V represents a large-scale effort, combining VQGAN and a 1.7 billion parameter denoising UNet to adapt to varying frame numbers via spatio-temporal blocks [\texttt{arXiv:2308.06571}].

\subsection{Controllability: Motion, Camera, and Compositional Control}

Controllability in video diffusion models has evolved from simple text-guided generation toward precise, decoupled manipulation of spatial composition, object trajectories, and camera movements. Early approaches focused on temporal consistency through zero-shot editing, such as Pix2Video, which uses self-attention feature injection to propagate changes from an anchor frame [\texttt{arXiv:2303.12688}], and vid2vid-zero, which employs a cross-frame modeling module to enable bi-directional temporal modeling at test time without video-specific training [\texttt{arXiv:2303.17599}]. While these methods are computationally efficient due to their training-free nature, they primarily focus on editing existing content rather than synthesizing complex, multi-modal motions from scratch.

A significant shift toward fine-grained control is seen in models that integrate explicit geometric or motion signals. VideoComposer introduces a Spatio-Temporal Condition (STC) encoder to unify spatial and temporal relations, utilizing motion vectors from compressed videos as explicit control signals [\texttt{arXiv:2306.02018}]. This move toward explicit guidance is further refined by DragNUWA, which addresses the limitation of existing works that focus on only one modality at a time by integrating text, image, and trajectory through an Adaptive Training (AT) strategy [\texttt{arXiv:2308.08089}]. However, DragNUWA is noted to be limited by its reliance on simple datasets and a lack of complex curve capabilities in open-domain settings [\texttt{arXiv:2308.08089}].

The decoupling of camera and object motion represents a more recent frontier in controllability. MotionCtrl provides a unified framework for independent control of camera and object motion using appearance-free conditions [\texttt{arXiv:2312.03641}]. Similarly, Direct-a-Video achieves decoupled control by utilizing spatial cross-attention modulation for object motion and new temporal cross-attention layers to interpret quantitative camera movement parameters [\texttt{arXiv:2402.03162}]. CameraCtrl further specializes this by offering a plug-and-play camera pose control module that can be trained on top of existing video diffusion models [\texttt{arXiv:2404.02101}].

Compositional control also extends to character pose and structural arrangement. Follow Your Pose enables pose-controllable generation by using a two-stage training scheme that leverages a zero-initialized convolutional encoder for pose information [\texttt{arXiv:2304.01186}]. For long-form content, NUWA-XL employs a "Diffusion over Diffusion" architecture, using a global model for keyframes and local models to fill content, which significantly improves efficiency—reducing inference time for 1024 frames from 7.55 minutes to 26 seconds [\texttt{arXiv:2303.12346}]. Finally, VideoTetris addresses the challenge of complex semantics by manipulating attention maps spatially and temporally, though it remains limited in handling videos with dynamic changes in object numbers [\texttt{arXiv:2406.04277}].

\subsection{Scaling and Efficiency: Transformers and Long-Form Video}

Recent advancements in video diffusion models have shifted focus from short-clip generation toward scaling temporal duration and improving computational efficiency through architectural innovations. A primary challenge in long-form synthesis is managing the quadratic complexity of temporal attention and maintaining consistency over extended sequences.

Several approaches address this by moving away from standard attention mechanisms toward more efficient architectures. Latte [\texttt{arXiv:2401.03048}] introduces a Latent Diffusion Transformer (LDT) that utilizes spatio-temporal tokens and decomposes spatial and temporal dimensions into four efficient variants to handle large token counts. Similarly, Open-Sora [\texttt{arXiv:2412.20404}] employs a Spatial-Temporal Diffusion Transformer (STDiT) with decoupled spatial and temporal attention, paired with a highly compressive 3D autoencoder to achieve high-fidelity 720p video. While these Transformer-based models scale via decomposition, other researchers seek to replace attention entirely to achieve linear scaling. For instance, [\texttt{arXiv:2403.07711}] replaces quadratic-complexity attention with bidirectional State-Space Models (SSMs), such as Mamba, as temporal feature extractors. This allows for linear-time memory consumption relative to sequence length, demonstrating superior efficiency for sequences up to 256 frames compared to traditional attention-based models.

Scaling the temporal dimension also involves novel training and inference paradigms. NUWA-XL [\texttt{arXiv:2303.12346}] utilizes a "Diffusion over Diffusion" architecture, employing a global model for keyframes and local models to recursively fill content. This coarse-to-fine approach enables parallel generation and significantly reduces inference time, decreasing the average duration for 1024 frames from 7.55 minutes to 26 seconds (a 94.26\% reduction). In contrast to these structural changes, FreeNoise [\texttt{arXiv:2310.15169}] offers a tuning-free paradigm that extends pretrained models via noise rescheduling and motion injection. FreeNoise achieves long-range correlation with only a 17\% increase in time cost, a significant improvement over the 255\% cost seen in previous state-of-the-art methods.

Further refinements focus on the granularity of temporal modeling. While CogVideoX [\texttt{arXiv:2408.06072}] leverages an expert transformer and a 3D VAE for improved compression and deep text-video modality fusion, [\texttt{arXiv:2410.03160}] introduces a Vectorized Timestep Approach (VTV). This allows each frame to follow an independent noise schedule, providing more precise temporal dependency modeling than uniform schedules. Additionally, while models like MotionCtrl [\texttt{arXiv:2312.03641}] focus on the flexibility of camera and object trajectories, the broader trend suggests that the integration of highly compressive 3D autoencoders [\texttt{arXiv:2408.06072}, \texttt{arXiv:2412.20404}] and efficient temporal extractors [\texttt{arXiv:2403.07711}] is essential for the next generation of long-form video synthesis.

\subsection{Evaluation, Benchmarking, and Physical Consistency}

As video generation models scale, evaluation methodologies have shifted from general quality assessment toward specialized frameworks that dissect compositionality and physical reasoning. Early efforts to address the lack of human-aligned metrics and the inability of existing tools to provide specific developmental insights led to the development of VBench (\texttt{arXiv:2311.17982}), which introduced a suite of 16 hierarchical and disentangled dimensions for video quality. This framework was subsequently expanded by VBench++ (\texttt{arXiv:2411.13503}), which increased versatility by supporting both text-to-video and image-to-video modalities and incorporating human preference alignment. While these suites provide a broad view of generative quality, T2V-CompBench (\texttt{arXiv:2407.14505}) offers a more granular focus on compositional text-to-video generation, utilizing 1400 prompts across seven categories to highlight the specific difficulties models face when adhering to complex, multi-object instructions.

A significant gap remains between visual fidelity and the adherence to physical laws. VideoPhy (\texttt{arXiv:2406.03520}) addresses this by introducing the VideoPhy benchmark and the VideoCon-Physics auto-evaluator to assess physical commonsense. The limitations of current models are stark; for instance, CogVideoX-5B adheres to both text prompts and physical laws in only 39.6\% of instances (\texttt{arXiv:2406.03520}). This failure is not merely a matter of prompt adherence but suggests a deeper lack of rule abstraction. Research using a 2D simulation testbed (\texttt{arXiv:2411.02385}) indicates that diffusion-based models often exhibit "case-based" generalization—mimicking training examples—rather than discovering or abstracting general physical rules, causing them to fail in out-of-distribution scenarios.

This distinction between visual mimicry and true physical understanding is further emphasized by the Physics-IQ benchmark (\texttt{arXiv:2501.09038}). By testing five specific domains—fluid dynamics, optics, solid mechanics, magnetism, and thermodynamics—Physics-IQ reveals that the physical understanding of prominent models, including Sora, Runway, and Stable Video Diffusion, is severely limited. Crucially, these findings suggest that a model's ability to generate visually realistic content is largely unrelated to its ability to respect physical principles (\texttt{arXiv:2501.09038}). Consequently, while large-scale models like HunyuanVideo (\texttt{arXiv:2412.03603}) continue to bridge the gap in generative capability through massive parameter counts, the literature suggests that scaling alone has not yet resolved the fundamental disconnect between visual realism and physical consistency.

\subsection{The Current Frontier: Advanced Foundation Models and Real-time Synthesis}

The landscape of video diffusion has shifted from basic temporal modeling toward large-scale foundation models that balance high-fidelity synthesis with computational efficiency. Recent advancements in Diffusion Transformer (DiT) architectures have enabled significant scaling; for instance, Step-Video-T2V utilizes a 30B parameter DiT with 3D full attention and Video-DPO to mitigate artifacts [\texttt{arXiv:2502.10248}], while the Wan suite offers scalable models ranging from 1.3B to 14B parameters [\texttt{arXiv:2503.20314}]. To address the high cost of these models, research has branched into two distinct directions: extreme compression for real-time performance and architectural refinements for temporal consistency.

A primary frontier is the pursuit of real-time synthesis through high-ratio latent compression. LTX-Video achieves this by integrating a Video-VAE with a denoising transformer, utilizing a 1:192 compression ratio (32 x 32 x 8 pixels per token) to generate 5 seconds of 24 fps video in just 2 seconds on an H100 GPU [\texttt{arXiv:2501.00103}]. However, this efficiency comes with a trade-off, as high compression inherently limits the representation of fine details [\texttt{arXiv:2501.00103}]. In contrast, models like UniVid focus on quality through a pyramid diffusion approach, employing temporal-pyramid cross-frame spatial-temporal attention and a dual-stream cross-attention mechanism to unify text-to-video and image-to-video tasks [\texttt{arXiv:2603.13739}].

Temporal consistency remains a central challenge, addressed through both architectural and training-free methods. RepVideo improves spatial appearance and temporal stability by accumulating features from neighboring layers to enrich the attention mechanism [\texttt{arXiv:2501.08994}]. Conversely, training-free approaches seek to enhance coherence without model retraining; for example, a motion consistency loss based on inter-frame feature correlation can improve motion control [\texttt{arXiv:2501.07563}], while DiffuseSlide utilizes noise re-injection and sliding window latent denoising to generate high frame rates [\texttt{arXiv:2506.01454}]. Notably, DiffuseSlide is limited by flickering and degradation in fast-motion, long-sequence scenarios [\texttt{arXiv:2506.01454}].

For long-horizon generation, researchers are moving beyond simple autoregressive methods. Context Forcing addresses the student-teacher mismatch in autoregressive generation by ensuring the teacher has access to full generation history [\texttt{arXiv:2602.06028}], while LongLive-RAG treats long video generation as a retrieval-augmented generation (RAG) problem, using a searchable history of latents to reduce identity drift [\texttt{arXiv:2606.02553}]. These developments are categorized by recent surveys into four efficiency paradigms: step distillation, efficient attention, model compression, and cache/trajectory optimization [\texttt{arXiv:2604.15911}]. While these methods mitigate the multiplicative costs of spatial and temporal dimensions, a critical weakness remains: the lack of inherently stable objectives for long-horizon generation, leaving the field reliant on "stability engineering" rather than unified algorithmic solutions [\texttt{arXiv:2604.15911}].

\subsection{Tensions and Open Problems}

A primary tension in the literature concerns the architectural strategy for maintaining temporal coherence versus achieving high-fidelity spatial detail. One school of thought advocates for the direct extension of high-resolution image priors, such as the approach in \texttt{arXiv:2312.04483}, which utilizes hierarchical spatio-temporal decoupling to leverage the robust spatial representations found in models like \texttt{arXiv:2307.01952}. Conversely, other researchers argue that such extensions are insufficient for complex motion, suggesting that video generation requires more explicit factorization of text and motion to avoid the semantic drift often seen in purely spatial-centric models, as argued in \texttt{arXiv:2311.10709}. This debate is further complicated by the struggle to balance training efficiency with generative quality; while \texttt{arXiv:2209.14792} demonstrates that video generation can be achieved without massive text-video datasets, this approach is critiqued by \texttt{arXiv:2302.07685}, which suggests that projecting diffusion into latent spaces provides a more mathematically robust framework for handling the high dimensionality of video compared to the more heuristic methods used in \texttt{arXiv:2211.11018} or \texttt{arXiv:2211.13221}.

Another significant disagreement centers on the methodology for long-form video synthesis and temporal continuity. Early attempts at long-duration generation, such as the time-agnostic VQGAN approach in \texttt{arXiv:2204.03638}, have been challenged by more recent hierarchical frameworks like \texttt{arXiv:2303.12346}, which posits that "diffusion over diffusion" is necessary to prevent the structural decay inherent in simpler autoregressive or VQGAN-based methods. This tension extends to the realm of video editing and zero-shot transformation. While some frameworks rely on cross-attention control to manipulate existing content, as seen in \texttt{arXiv:2303.04761}, others argue that these methods lack the necessary temporal stability for video-to-video tasks. For instance, \texttt{arXiv:2308.09592} critiques several foundational editing and generation models—including \texttt{arXiv:2209.14792}, \texttt{arXiv:2302.01329}, and \texttt{arXiv:2302.03011}—arguing that they fail to maintain strict consistency-aware constraints, leading to flickering or semantic loss during the transformation process.

\subsection{Outlook}

The field is currently transitioning from simple text-to-video synthesis toward highly controllable, long-context, and memory-augmented generation. The frontier lies in moving beyond stochastic sampling toward deterministic motion alignment and sophisticated camera control, as suggested by recent advancements in motion-centric representation alignment (\texttt{arXiv:2510.19022}) and memory-augmented camera systems (\texttt{arXiv:2603.26193}). As we move toward 2026, the integration of retrieval-augmented frameworks (\texttt{arXiv:2606.02553}) and pyramid-based diffusion architectures (\texttt{arXiv:2603.13739}) promises to bridge the gap between short-clip generation and the creation of coherent, cinematic-length narratives.

\section*{References}
\begin{enumerate}\small
  \item Diffractive Scattering (2004). \texttt{arXiv:hep-ex/0410002}
  \item Inclusive and exclusive diffraction: overview of the experimental
  presentations at Diffraction2004 (2005). \texttt{arXiv:hep-ex/0501008}
  \item UCF101: A Dataset of 101 Human Actions Classes From Videos in The Wild (2012). \texttt{arXiv:1212.0402}
  \item Auto-Encoding Variational Bayes (2013). \texttt{arXiv:1312.6114}
  \item Generative Adversarial Networks (2014). \texttt{arXiv:1406.2661}
  \item Deep Unsupervised Learning using Nonequilibrium Thermodynamics (2015). \texttt{arXiv:1503.03585}
  \item U-Net: Convolutional Networks for Biomedical Image Segmentation (2015). \texttt{arXiv:1505.04597}
  \item Generating Videos with Scene Dynamics (2016). \texttt{arXiv:1609.02612}
  \item Temporal Generative Adversarial Nets with Singular Value Clipping (2016). \texttt{arXiv:1611.06624}
  \item Sync-DRAW: Automatic Video Generation using Deep Recurrent Attentive
  Architectures (2016). \texttt{arXiv:1611.10314}
  \item Attention Is All You Need (2017). \texttt{arXiv:1706.03762}
  \item GANs Trained by a Two Time-Scale Update Rule Converge to a Local Nash
  Equilibrium (2017). \texttt{arXiv:1706.08500}
  \item Decoupled Weight Decay Regularization (2017). \texttt{arXiv:1711.05101}
  \item Neural Discrete Representation Learning (2017). \texttt{arXiv:1711.00937}
  \item MoCoGAN: Decomposing Motion and Content for Video Generation (2017). \texttt{arXiv:1707.04993}
  \item Quo Vadis, Action Recognition? A New Model and the Kinetics Dataset (2017). \texttt{arXiv:1705.07750}
  \item The 2017 DAVIS Challenge on Video Object Segmentation (2017). \texttt{arXiv:1704.00675}
  \item Video Generation From Text (2017). \texttt{arXiv:1710.00421}
  \item Diffractive Scattering: Problems in Theory and Praxis (2017). \texttt{arXiv:1711.00226}
  \item A Style-Based Generator Architecture for Generative Adversarial Networks (2018). \texttt{arXiv:1812.04948}
  \item The Unreasonable Effectiveness of Deep Features as a Perceptual Metric (2018). \texttt{arXiv:1801.03924}
  \item Towards Accurate Generative Models of Video: A New Metric \& Challenges (2018). \texttt{arXiv:1812.01717}
  \item Video-to-Video Synthesis (2018). \texttt{arXiv:1808.06601}
  \item Train Sparsely, Generate Densely: Memory-efficient Unsupervised Training
  of High-resolution Temporal GAN (2018). \texttt{arXiv:1811.09245}
  \item Learning to Forecast and Refine Residual Motion for Image-to-Video
  Generation (2018). \texttt{arXiv:1807.09951}
  \item Generative Modeling by Estimating Gradients of the Data Distribution (2019). \texttt{arXiv:1907.05600}
  \item Analyzing and Improving the Image Quality of StyleGAN (2019). \texttt{arXiv:1912.04958}
  \item Exploring the Limits of Transfer Learning with a Unified Text-to-Text
  Transformer (2019). \texttt{arXiv:1910.10683}
  \item Adversarial Video Generation on Complex Datasets (2019). \texttt{arXiv:1907.06571}
  \item Markov Decision Process for Video Generation (2019). \texttt{arXiv:1909.12400}
  \item Few-shot Video-to-Video Synthesis (2019). \texttt{arXiv:1910.12713}
  \item Denoising Diffusion Probabilistic Models (2020). \texttt{arXiv:2006.11239}
  \item Denoising Diffusion Implicit Models (2020). \texttt{arXiv:2010.02502}
  \item Score-Based Generative Modeling through Stochastic Differential
  Equations (2020). \texttt{arXiv:2011.13456}
  \item An Image is Worth 16x16 Words: Transformers for Image Recognition at
  Scale (2020). \texttt{arXiv:2010.11929}
  \item Taming Transformers for High-Resolution Image Synthesis (2020). \texttt{arXiv:2012.09841}
  \item RAFT: Recurrent All-Pairs Field Transforms for Optical Flow (2020). \texttt{arXiv:2003.12039}
  \item Latent Video Transformer (2020). \texttt{arXiv:2006.10704}
  \item First Order Motion Model for Image Animation (2020). \texttt{arXiv:2003.00196}
  \item Interactive Video Stylization Using Few-Shot Patch-Based Training (2020). \texttt{arXiv:2004.14489}
  \item Temporal Shift GAN for Large Scale Video Generation (2020). \texttt{arXiv:2004.01823}
  \item Learning Transferable Visual Models From Natural Language Supervision (2021). \texttt{arXiv:2103.00020}
  \item High-Resolution Image Synthesis with Latent Diffusion Models (2021). \texttt{arXiv:2112.10752}
  \item Diffusion Models Beat GANs on Image Synthesis (2021). \texttt{arXiv:2105.05233}
  \item Zero-Shot Text-to-Image Generation (2021). \texttt{arXiv:2102.12092}
  \item GLIDE: Towards Photorealistic Image Generation and Editing with
  Text-Guided Diffusion Models (2021). \texttt{arXiv:2112.10741}
  \item Improved Denoising Diffusion Probabilistic Models (2021). \texttt{arXiv:2102.09672}
  \item Frozen in Time: A Joint Video and Image Encoder for End-to-End Retrieval (2021). \texttt{arXiv:2104.00650}
  \item Cascaded Diffusion Models for High Fidelity Image Generation (2021). \texttt{arXiv:2106.15282}
  \item VideoGPT: Video Generation using VQ-VAE and Transformers (2021). \texttt{arXiv:2104.10157}
  \item SDEdit: Guided Image Synthesis and Editing with Stochastic Differential
  Equations (2021). \texttt{arXiv:2108.01073}
  \item N\textbackslash{}"UWA: Visual Synthesis Pre-training for Neural visUal World creAtion (2021). \texttt{arXiv:2111.12417}
  \item GODIVA: Generating Open-DomaIn Videos from nAtural Descriptions (2021). \texttt{arXiv:2104.14806}
  \item LoRA: Low-Rank Adaptation of Large Language Models (2021). \texttt{arXiv:2106.09685}
  \item A Good Image Generator Is What You Need for High-Resolution Video
  Synthesis (2021). \texttt{arXiv:2104.15069}
  \item StyleGAN-V: A Continuous Video Generator with the Price, Image Quality
  and Perks of StyleGAN2 (2021). \texttt{arXiv:2112.14683}
  \item Advancing High-Resolution Video-Language Representation with Large-Scale
  Video Transcriptions (2021). \texttt{arXiv:2111.10337}
  \item StyleVideoGAN: A Temporal Generative Model using a Pretrained StyleGAN (2021). \texttt{arXiv:2107.07224}
  \item Make-A-Video: Text-to-Video Generation without Text-Video Data (2022). \texttt{arXiv:2209.14792}
  \item Video Diffusion Models (2022). \texttt{arXiv:2204.03458}
  \item Imagen Video: High Definition Video Generation with Diffusion Models (2022). \texttt{arXiv:2210.02303}
  \item Photorealistic Text-to-Image Diffusion Models with Deep Language
  Understanding (2022). \texttt{arXiv:2205.11487}
  \item Hierarchical Text-Conditional Image Generation with CLIP Latents (2022). \texttt{arXiv:2204.06125}
  \item Classifier-Free Diffusion Guidance (2022). \texttt{arXiv:2207.12598}
  \item CogVideo: Large-scale Pretraining for Text-to-Video Generation via
  Transformers (2022). \texttt{arXiv:2205.15868}
  \item Tune-A-Video: One-Shot Tuning of Image Diffusion Models for
  Text-to-Video Generation (2022). \texttt{arXiv:2212.11565}
  \item MagicVideo: Efficient Video Generation With Latent Diffusion Models (2022). \texttt{arXiv:2211.11018}
  \item Latent Video Diffusion Models for High-Fidelity Long Video Generation (2022). \texttt{arXiv:2211.13221}
  \item Scalable Diffusion Models with Transformers (2022). \texttt{arXiv:2212.09748}
  \item DreamBooth: Fine Tuning Text-to-Image Diffusion Models for
  Subject-Driven Generation (2022). \texttt{arXiv:2208.12242}
  \item Prompt-to-Prompt Image Editing with Cross Attention Control (2022). \texttt{arXiv:2208.01626}
  \item Long Video Generation with Time-Agnostic VQGAN and Time-Sensitive
  Transformer (2022). \texttt{arXiv:2204.03638}
  \item LAION-5B: An open large-scale dataset for training next generation
  image-text models (2022). \texttt{arXiv:2210.08402}
  \item Generating Videos with Dynamics-aware Implicit Generative Adversarial
  Networks (2022). \texttt{arXiv:2202.10571}
  \item Progressive Distillation for Fast Sampling of Diffusion Models (2022). \texttt{arXiv:2202.00512}
  \item Flexible Diffusion Modeling of Long Videos (2022). \texttt{arXiv:2205.11495}
  \item Phenaki: Variable Length Video Generation From Open Domain Textual
  Description (2022). \texttt{arXiv:2210.02399}
  \item An Image is Worth One Word: Personalizing Text-to-Image Generation using
  Textual Inversion (2022). \texttt{arXiv:2208.01618}
  \item Elucidating the Design Space of Diffusion-Based Generative Models (2022). \texttt{arXiv:2206.00364}
  \item Text2LIVE: Text-Driven Layered Image and Video Editing (2022). \texttt{arXiv:2204.02491}
  \item MCVD: Masked Conditional Video Diffusion for Prediction, Generation, and
  Interpolation (2022). \texttt{arXiv:2205.09853}
  \item eDiff-I: Text-to-Image Diffusion Models with an Ensemble of Expert
  Denoisers (2022). \texttt{arXiv:2211.01324}
  \item Diffusion Probabilistic Modeling for Video Generation (2022). \texttt{arXiv:2203.09481}
  \item InstructPix2Pix: Learning to Follow Image Editing Instructions (2022). \texttt{arXiv:2211.09800}
  \item Plug-and-Play Diffusion Features for Text-Driven Image-to-Image
  Translation (2022). \texttt{arXiv:2211.12572}
  \item Generating Long Videos of Dynamic Scenes (2022). \texttt{arXiv:2206.03429}
  \item MAGVIT: Masked Generative Video Transformer (2022). \texttt{arXiv:2212.05199}
  \item New features in the differential cross sections measured at the LHC (2022). \texttt{arXiv:2210.02970}
  \item Diffusion Models for Video Prediction and Infilling (2022). \texttt{arXiv:2206.07696}
  \item SinFusion: Training Diffusion Models on a Single Image or Video (2022). \texttt{arXiv:2211.11743}
  \item Align your Latents: High-Resolution Video Synthesis with Latent
  Diffusion Models (2023). \texttt{arXiv:2304.08818}
  \item Hierarchical Spatio-temporal Decoupling for Text-to-Video Generation (2023). \texttt{arXiv:2312.04483}
  \item Stable Video Diffusion: Scaling Latent Video Diffusion Models to Large
  Datasets (2023). \texttt{arXiv:2311.15127}
  \item VideoGen: A Reference-Guided Latent Diffusion Approach for High
  Definition Text-to-Video Generation (2023). \texttt{arXiv:2309.00398}
  \item MoVideo: Motion-Aware Video Generation with Diffusion Models (2023). \texttt{arXiv:2311.11325}
  \item LAVIE: High-Quality Video Generation with Cascaded Latent Diffusion
  Models (2023). \texttt{arXiv:2309.15103}
  \item Swap Attention in Spatiotemporal Diffusions for Text-to-Video Generation (2023). \texttt{arXiv:2305.10874}
  \item Video Probabilistic Diffusion Models in Projected Latent Space (2023). \texttt{arXiv:2302.07685}
  \item Adding Conditional Control to Text-to-Image Diffusion Models (2023). \texttt{arXiv:2302.05543}
  \item AnimateDiff: Animate Your Personalized Text-to-Image Diffusion Models
  without Specific Tuning (2023). \texttt{arXiv:2307.04725}
  \item Structure and Content-Guided Video Synthesis with Diffusion Models (2023). \texttt{arXiv:2302.03011}
  \item VBench: Comprehensive Benchmark Suite for Video Generative Models (2023). \texttt{arXiv:2311.17982}
  \item ModelScope Text-to-Video Technical Report (2023). \texttt{arXiv:2308.06571}
  \item Text2Video-Zero: Text-to-Image Diffusion Models are Zero-Shot Video
  Generators (2023). \texttt{arXiv:2303.13439}
  \item VideoComposer: Compositional Video Synthesis with Motion Controllability (2023). \texttt{arXiv:2306.02018}
  \item VideoCrafter1: Open Diffusion Models for High-Quality Video Generation (2023). \texttt{arXiv:2310.19512}
  \item SDXL: Improving Latent Diffusion Models for High-Resolution Image
  Synthesis (2023). \texttt{arXiv:2307.01952}
  \item T2I-Adapter: Learning Adapters to Dig out More Controllable Ability for
  Text-to-Image Diffusion Models (2023). \texttt{arXiv:2302.08453}
  \item VideoFusion: Decomposed Diffusion Models for High-Quality Video
  Generation (2023). \texttt{arXiv:2303.08320}
  \item Preserve Your Own Correlation: A Noise Prior for Video Diffusion Models (2023). \texttt{arXiv:2305.10474}
  \item FateZero: Fusing Attentions for Zero-shot Text-based Video Editing (2023). \texttt{arXiv:2303.09535}
  \item Show-1: Marrying Pixel and Latent Diffusion Models for Text-to-Video Generation (2023). \texttt{arXiv:2309.15818}
  \item Dreamix: Video Diffusion Models are General Video Editors (2023). \texttt{arXiv:2302.01329}
  \item MotionCtrl: A Unified and Flexible Motion Controller for Video
  Generation (2023). \texttt{arXiv:2312.03641}
  \item DynamiCrafter: Animating Open-domain Images with Video Diffusion Priors (2023). \texttt{arXiv:2310.12190}
  \item DragNUWA: Fine-grained Control in Video Generation by Integrating Text,
  Image, and Trajectory (2023). \texttt{arXiv:2308.08089}
  \item I2VGen-XL: High-Quality Image-to-Video Synthesis via Cascaded Diffusion
  Models (2023). \texttt{arXiv:2311.04145}
  \item ControlVideo: Training-free Controllable Text-to-Video Generation (2023). \texttt{arXiv:2305.13077}
  \item NUWA-XL: Diffusion over Diffusion for eXtremely Long Video Generation (2023). \texttt{arXiv:2303.12346}
  \item Pix2Video: Video Editing using Image Diffusion (2023). \texttt{arXiv:2303.12688}
  \item InternVid: A Large-scale Video-Text Dataset for Multimodal Understanding
  and Generation (2023). \texttt{arXiv:2307.06942}
  \item Video-P2P: Video Editing with Cross-attention Control (2023). \texttt{arXiv:2303.04761}
  \item MVBench: A Comprehensive Multi-modal Video Understanding Benchmark (2023). \texttt{arXiv:2311.17005}
  \item Emu Video: Factorizing Text-to-Video Generation by Explicit Image
  Conditioning (2023). \texttt{arXiv:2311.10709}
  \item TokenFlow: Consistent Diffusion Features for Consistent Video Editing (2023). \texttt{arXiv:2307.10373}
  \item EvalCrafter: Benchmarking and Evaluating Large Video Generation Models (2023). \texttt{arXiv:2310.11440}
  \item PixArt-\$\textbackslash{}alpha\$: Fast Training of Diffusion Transformer for
  Photorealistic Text-to-Image Synthesis (2023). \texttt{arXiv:2310.00426}
  \item Photorealistic Video Generation with Diffusion Models (2023). \texttt{arXiv:2312.06662}
  \item Gen-L-Video: Multi-Text to Long Video Generation via Temporal
  Co-Denoising (2023). \texttt{arXiv:2305.18264}
  \item SEINE: Short-to-Long Video Diffusion Model for Generative Transition and
  Prediction (2023). \texttt{arXiv:2310.20700}
  \item Zero-Shot Video Editing Using Off-The-Shelf Image Diffusion Models (2023). \texttt{arXiv:2303.17599}
  \item Rerender A Video: Zero-Shot Text-Guided Video-to-Video Translation (2023). \texttt{arXiv:2306.07954}
  \item FreeNoise: Tuning-Free Longer Video Diffusion via Noise Rescheduling (2023). \texttt{arXiv:2310.15169}
  \item Follow Your Pose: Pose-Guided Text-to-Video Generation using Pose-Free
  Videos (2023). \texttt{arXiv:2304.01186}
  \item Animate Anyone: Consistent and Controllable Image-to-Video Synthesis for
  Character Animation (2023). \texttt{arXiv:2311.17117}
  \item Latent-Shift: Latent Diffusion with Temporal Shift for Efficient
  Text-to-Video Generation (2023). \texttt{arXiv:2304.08477}
  \item Control-A-Video: Controllable Text-to-Video Diffusion Models with Motion
  Prior and Reward Feedback Learning (2023). \texttt{arXiv:2305.13840}
  \item MotionDirector: Motion Customization of Text-to-Video Diffusion Models (2023). \texttt{arXiv:2310.08465}
  \item Motion-Conditioned Diffusion Model for Controllable Video Synthesis (2023). \texttt{arXiv:2304.14404}
  \item Make Pixels Dance: High-Dynamic Video Generation (2023). \texttt{arXiv:2311.10982}
  \item Make-Your-Video: Customized Video Generation Using Textual and
  Structural Guidance (2023). \texttt{arXiv:2306.00943}
  \item MagicAnimate: Temporally Consistent Human Image Animation using
  Diffusion Model (2023). \texttt{arXiv:2311.16498}
  \item StableVideo: Text-driven Consistency-aware Diffusion Video Editing (2023). \texttt{arXiv:2308.09592}
  \item VDT: General-purpose Video Diffusion Transformers via Mask Modeling (2023). \texttt{arXiv:2305.13311}
  \item LAMP: Learn A Motion Pattern for Few-Shot-Based Video Generation (2023). \texttt{arXiv:2310.10769}
  \item DreamPose: Fashion Image-to-Video Synthesis via Stable Diffusion (2023). \texttt{arXiv:2304.06025}
  \item DreamVideo: Composing Your Dream Videos with Customized Subject and
  Motion (2023). \texttt{arXiv:2312.04433}
  \item VideoDirectorGPT: Consistent Multi-scene Video Generation via LLM-Guided
  Planning (2023). \texttt{arXiv:2309.15091}
  \item Animate-A-Story: Storytelling with Retrieval-Augmented Video Generation (2023). \texttt{arXiv:2307.06940}
  \item FETV: A Benchmark for Fine-Grained Evaluation of Open-Domain
  Text-to-Video Generation (2023). \texttt{arXiv:2311.01813}
  \item Reuse and Diffuse: Iterative Denoising for Text-to-Video Generation (2023). \texttt{arXiv:2309.03549}
  \item Video Diffusion Models: A Survey (2024). \texttt{arXiv:2405.03150}
  \item Grid Diffusion Models for Text-to-Video Generation (2024). \texttt{arXiv:2404.00234}
  \item Redefining Temporal Modeling in Video Diffusion: The Vectorized Timestep
  Approach (2024). \texttt{arXiv:2410.03160}
  \item VideoTetris: Towards Compositional Text-to-Video Generation (2024). \texttt{arXiv:2406.04277}
  \item Lumiere: A Space-Time Diffusion Model for Video Generation (2024). \texttt{arXiv:2401.12945}
  \item ExVideo: Extending Video Diffusion Models via Parameter-Efficient
  Post-Tuning (2024). \texttt{arXiv:2406.14130}
  \item Efficient Video Diffusion Models via Content-Frame Motion-Latent
  Decomposition (2024). \texttt{arXiv:2403.14148}
  \item SSM Meets Video Diffusion Models: Efficient Long-Term Video Generation
  with Structured State Spaces (2024). \texttt{arXiv:2403.07711}
  \item CogVideoX: Text-to-Video Diffusion Models with An Expert Transformer (2024). \texttt{arXiv:2408.06072}
  \item HunyuanVideo: A Systematic Framework For Large Video Generative Models (2024). \texttt{arXiv:2412.03603}
  \item VideoCrafter2: Overcoming Data Limitations for High-Quality Video
  Diffusion Models (2024). \texttt{arXiv:2401.09047}
  \item Open-Sora: Democratizing Efficient Video Production for All (2024). \texttt{arXiv:2412.20404}
  \item Scaling Rectified Flow Transformers for High-Resolution Image Synthesis (2024). \texttt{arXiv:2403.03206}
  \item Movie Gen: A Cast of Media Foundation Models (2024). \texttt{arXiv:2410.13720}
  \item Panda-70M: Captioning 70M Videos with Multiple Cross-Modality Teachers (2024). \texttt{arXiv:2402.19479}
  \item Latte: Latent Diffusion Transformer for Video Generation (2024). \texttt{arXiv:2401.03048}
  \item CameraCtrl: Enabling Camera Control for Text-to-Video Generation (2024). \texttt{arXiv:2404.02101}
  \item VideoPhy: Evaluating Physical Commonsense for Video Generation (2024). \texttt{arXiv:2406.03520}
  \item StreamingT2V: Consistent, Dynamic, and Extendable Long Video Generation
  from Text (2024). \texttt{arXiv:2403.14773}
  \item Direct-a-Video: Customized Video Generation with User-Directed Camera
  Movement and Object Motion (2024). \texttt{arXiv:2402.03162}
  \item Towards World Simulator: Crafting Physical Commonsense-Based Benchmark
  for Video Generation (2024). \texttt{arXiv:2410.05363}
  \item Pyramidal Flow Matching for Efficient Video Generative Modeling (2024). \texttt{arXiv:2410.05954}
  \item Open-Sora Plan: Open-Source Large Video Generation Model (2024). \texttt{arXiv:2412.00131}
  \item Motion-I2V: Consistent and Controllable Image-to-Video Generation with
  Explicit Motion Modeling (2024). \texttt{arXiv:2401.15977}
  \item VBench++: Comprehensive and Versatile Benchmark Suite for Video
  Generative Models (2024). \texttt{arXiv:2411.13503}
  \item T2V-CompBench: A Comprehensive Benchmark for Compositional Text-to-video
  Generation (2024). \texttt{arXiv:2407.14505}
  \item How Far is Video Generation from World Model: A Physical Law Perspective (2024). \texttt{arXiv:2411.02385}
  \item Current Status of the Odderon (2024). \texttt{arXiv:2408.01990}
  \item VideoScore: Building Automatic Metrics to Simulate Fine-grained Human
  Feedback for Video Generation (2024). \texttt{arXiv:2406.15252}
  \item Survey of Video Diffusion Models: Foundations, Implementations, and Applications (2025). \texttt{arXiv:2504.16081}
  \item From Image to Video: An Empirical Study of Diffusion Representations (2025). \texttt{arXiv:2502.07001}
  \item RepVideo: Rethinking Cross-Layer Representation for Video Generation (2025). \texttt{arXiv:2501.08994}
  \item MoAlign: Motion-Centric Representation Alignment for Video Diffusion Models (2025). \texttt{arXiv:2510.19022}
  \item Emergent Temporal Correspondences from Video Diffusion Transformers (2025). \texttt{arXiv:2506.17220}
  \item DiffuseSlide: Training-Free High Frame Rate Video Generation Diffusion (2025). \texttt{arXiv:2506.01454}
  \item Temporal In-Context Fine-Tuning with Temporal Reasoning for Versatile Control of Video Diffusion Models (2025). \texttt{arXiv:2506.00996}
  \item Training-Free Motion-Guided Video Generation with Enhanced Temporal
  Consistency Using Motion Consistency Loss (2025). \texttt{arXiv:2501.07563}
  \item Wan: Open and Advanced Large-Scale Video Generative Models (2025). \texttt{arXiv:2503.20314}
  \item LTX-Video: Realtime Video Latent Diffusion (2025). \texttt{arXiv:2501.00103}
  \item Do generative video models understand physical principles? (2025). \texttt{arXiv:2501.09038}
  \item Step-Video-T2V Technical Report: The Practice, Challenges, and Future of
  Video Foundation Model (2025). \texttt{arXiv:2502.10248}
  \item Efficient Video Diffusion Models: Advancements and Challenges (2026). \texttt{arXiv:2604.15911}
  \item Image-to-Video Diffusion: From Foundations to Open Frontiers (2026). \texttt{arXiv:2605.17248}
  \item UniVid: Pyramid Diffusion Model for High Quality Video Generation (2026). \texttt{arXiv:2603.13739}
  \item MoRight: Motion Control Done Right (2026). \texttt{arXiv:2604.07348}
  \item V2V-Bench: A Comprehensive Benchmark for Video-to-Video Generation Evaluation (2026). \texttt{arXiv:2606.05665}
  \item Context Forcing: Consistent Autoregressive Video Generation with Long Context (2026). \texttt{arXiv:2602.06028}
  \item LongLive-RAG: A General Retrieval-Augmented Framework for Long Video Generation (2026). \texttt{arXiv:2606.02553}
  \item MemCam: Memory-Augmented Camera Control for Consistent Video Generation (2026). \texttt{arXiv:2603.26193}
\end{enumerate}

\end{document}
